\section{Binary merit factors and Morse shifts}\label{sec:merit-morse}
We now derive the autocorrelation and symbolic-dynamical consequences of
the maximum estimate. For a binary word
$A=(a_0,\ldots,a_{N-1})\in\{-1,1\}^N$, recall
\[
 P_A(z)=\sum_{j=0}^{N-1}a_jz^j,
 \qquad C_u(A)=\sum_{j=0}^{N-1-u}a_ja_{j+u}\quad(1\le u<N).
\]
For $N\ge2$, its merit factor is
$F(A)=N^2/(2\sum_{u=1}^{N-1}C_u(A)^2)$. The denominator is positive, since
$C_{N-1}(A)=a_0a_{N-1}\in\{-1,1\}$. Expanding $|P_A|^2$ on the unit
circle and applying Parseval with probability Haar measure gives
\begin{equation}\label{eq:merit-parseval}
 \norm{P_A}_4^4=N^2+2\sum_{u=1}^{N-1}C_u(A)^2,
 \qquad
 F(A)=\frac{N^2}{\norm{P_A}_4^4-N^2}.
\end{equation}
In the notation of Downarowicz--Lacroix,
$M_A=N^{-2}\sum_{u=1}^{N-1}C_u(A)^2$ and $F(A)=1/(2M_A)$
\cite[Definition~2 and Lemma~0]{DownarowiczLacroix1998}.

\begin{corollary}[All-length unbounded binary merit factors]\label{cor:merit}
For $N\ge2$, let
\[
 \mathcal F_N=\max_{A\in\{-1,1\}^N}F(A).
\]
Then
\[
 \lim_{N\to\infty}\mathcal F_N=+\infty,
\]
where the limit runs through all integer lengths. Equivalently, one can
choose a word $A_N\in\{-1,1\}^N$ for each $N\ge2$ such that
\[
 \sum_{u=1}^{N-1}C_u(A_N)^2=o(N^2).
\]
\end{corollary}
\begin{proof}
Choose $A_N$ so that $P_{A_N}$ attains the minimum defining $m_N$.
By \eqref{eq:merit-parseval}, Parseval, and the elementary fourth-moment
bound,
\[
 0<\frac1{F(A_N)}
 =\frac{\norm{P_{A_N}}_4^4}{N^2}-1
 \le\frac{\norm{P_{A_N}}_\infty^2\norm{P_{A_N}}_2^2}{N^2}-1
 =m_N^2-1\longrightarrow0.
\]
The last limit is Theorem~\ref{thm:main}. Thus $F(A_N)\to\infty$ and
$\mathcal F_N\ge F(A_N)$ has the same conclusion. The equivalence with
the autocorrelation formulation follows directly from the definition of
$F(A)$.
\end{proof}

This corollary refutes the bounded-merit-factor conjecture while retaining
the all-length quantifier of Theorem~\ref{thm:main}. It supplies no
quantitative rate of divergence.

For completeness, we specify the symbolic class in the next consequence.
For binary words $B$ and $C$ of lengths $b$ and $c$, respectively, their
product is defined by
\[
 (B\times C)(s+bt)=B(s)C(t)\qquad(0\le s<b,\ 0\le t<c).
\]
If each $B_j$ begins with $1$ and has length at least two, the products
$B_1\times\cdots\times B_j$ have a one-sided coordinatewise limit.
This limit admits a two-sided extension with the same finite words
\cite[Definition~3 and the following paragraph]{DownarowiczLacroix1998}.
The orbit closure $X\subset\{-1,1\}^{\Z}$ of such an extension, under
$Sx(n)=x(n+1)$, is a \emph{binary Morse shift}. This is the discrete-time
system called a Morse flow in that paper. Simple spectrum means that the
Koopman operator has a cyclic vector: the linear span of its integer iterates is
dense in the underlying $L^2$ space.

\begin{corollary}[Binary Morse spectral consequence]\label{cor:morse}
There is a binary Morse shift $(X,S)$ with a unique $S$-invariant Borel
probability measure $\mu$ for which the Koopman operator
$U_Sg=g\circ S$ on $L^2(X,\mu)$ has simple spectrum. For the
zero-coordinate function $f(x)=x(0)$, let $\sigma_f$ be the spectral
measure characterized by
\[
 \widehat{\sigma_f}(n):=\int_{\T}\e(nt)\,d\sigma_f(t)
 =\int_X f(x)f(S^n x)\,d\mu(x)\qquad(n\in\Z).
\]
Then $d\sigma_f=h\,dt$ for some $h\in L^2(\T)$ with $h\ge0$ almost
everywhere and $\int_{\T}h\,dt=1$. In particular, the absolutely
continuous spectral component is nonzero.
\end{corollary}
\begin{proof}
Corollary~\ref{cor:merit} supplies nontrivial finite binary words with
arbitrarily large merit factors. The forward construction in
\cite[Theorem~2 and its proof]{DownarowiczLacroix1998} therefore gives a
binary Morse shift with $\sum_{n=1}^{\infty}|\widehat{\sigma_f}(n)|^2<\infty$.
Its defining words have frequencies of both signs tending to $1/2$;
their Fact~1 gives unique ergodicity, and their Fact~2 gives simple
spectrum for that invariant measure.
Conjugate symmetry extends the square summability to $\Z$. Fourier
Plancherel and uniqueness of Fourier coefficients then give
$\sigma_f=h\,dt$ with $h\in L^2(\T)$. Positivity of $\sigma_f$ gives
$h\ge0$ almost everywhere, and
$\int h\,dt=\sigma_f(\T)=\norm{f}_2^2=1$.
\end{proof}

Here simplicity concerns the full Koopman operator, whereas absolute
continuity is asserted for the cyclic subspace generated by $f$; the
remaining spectral type is not identified. A different realization is
given in the companion manuscript \cite[Theorem~1.1]{OpenAIThreeTorus2026}:
it constructs a smooth volume-preserving diffeomorphism of $\T^3$ whose
Koopman operator has simple Lebesgue spectrum on the entire mean-zero
$L^2$ space. The present consequence instead concerns the constrained
class of binary Morse shifts.
